\documentclass{article}
\usepackage{graphicx} % Required for inserting images
\usepackage{caption} % Required for \captionof
\usepackage{amsmath}
\usepackage{float}
\usepackage{comment}
\usepackage[T1]{fontenc}
\usepackage[utf8]{inputenc}
\usepackage{lmodern}
\usepackage{microtype}
\usepackage{xcolor}
\usepackage{framed}
\renewenvironment{leftbar}{%
  \def\FrameCommand{\textcolor{orange}{\vrule width 3pt} \hspace{10pt}}%
  \MakeFramed {\advance\hsize-\width \FrameRestore}}%
 {\endMakeFramed}

\title{Mechanikai mérés}
\author{Kern Luca, Vincze Csongor}
\date{2026 Április 22}

\begin{document}

\maketitle

\section{10. Feladat: Vizsgálja meg a lebegés jelenségét!}

$d_{\text{mágnesek}_{\text{max}}} = 2,1 \text{cm}$
%mérőrúd!
$A_{\text{max}} = \text{mm}$
$f = \text{Hz}$    %f0-0,1Hz

\begin{figure}[H]
    \centering
    \includegraphics[width=4cm]{../images/10_feladat1.png}
    \caption*{Lebegés}
    \label{lebeges}
\end{figure}

\end{document}